\usepackage{hyperref,cleveref}
\usepackage{environ}
\NewEnviron{killcontents}{}
\usepackage{todonotes}
\presetkeys%
{todonotes}%
{inline,backgroundcolor=yellow}{}
\usepackage{graphicx}
\setlength{\marginparwidth}{2cm}
\usepackage{amsmath, amssymb, amsthm}
\usepackage{pifont}
\usepackage{bbold}
\usepackage{float}
\usepackage{mathtools}
\usepackage{ebproof}
\usepackage{braket}
\usepackage{comment}
\usepackage{xcolor}
\usepackage{float}
\usepackage{mfirstuc}
\usepackage{appendix}
\usepackage{thmtools, thm-restate}
\usepackage{xparse}
\usepackage[ruled, lined, noend]{algorithm2e}
\usepackage{array}
\usepackage{ebproof}
\usepackage{stmaryrd}

\hypersetup{colorlinks,linkcolor={blue},citecolor={blue}}
%%%%%%%%%%%%%%%%%%%%%%%%%%
%       MACROS
%%%%%%%%%%%%%%%%%%%%%%%%%%

%Prooftrees
\newcommand{\pt}[1]{
  \begin{prooftree}[center = false] #1
\end{prooftree}}
\newcommand{\ptsc}[1]{
  \begin{prooftree}[center=false, right label template = \small\inserttext,
    template = \small$\inserttext$] #1
\end{prooftree}}
\ebproofnewstyle{small}{
  rule margin = 1.5ex,
  separation = 0.5em,
  right label template = \footnotesize\inserttext,
  left label template = \footnotesize\inserttext,
  template = \footnotesize$\inserttext$ 
}

%Base types in the grammar
\newcommand{\cvar}{C}
\newcommand{\tvar}{T}
\newcommand{\bqvar}{B_Q}
\newcommand{\bcvar}{B_C}
\newcommand{\bvar}{B}
\newcommand{\qvar}{Q}
\newcommand\bkvar{\bvar^{\ket{}}}
\newcommand{\bkvarb}{{\bkvar}}
\newcommand{\qbit}{\mathrm{Qbit}}
\newcommand{\nonlinfunc}{\Rightarrow}
\newcommand{\linfunc}{\multimap}
\newcommand{\anyfunc}{\rightarrowtail}

%Basic type sets
\newcommand{\bset}{\mathbb{B}}
\newcommand{\bcset}{\mathbb{B}_C}
\newcommand{\bqset}{\mathbb{B}_Q}
\newcommand{\unit}{\mathbb 1}
\newcommand{\nat}{\mathtt{nat}}
\newcommand{\typelist}[1]{\mathtt{list}(#1)}
%Constructors
\newcommand{\cons}[1]{\mathrm{Cons}(#1)}
\newcommand{\sign}[2]{\ifx&#1&
  \else
  #1 \to
  \fi
#2}
\newcommand{\signdefault}{\sign {\bkvar_1, \dots, \bkvar_n} \bvar}
\newcommand{\vunit}{()}
\newcommand{\nil}{[~]}

%Terms
\newcommand{\textqcase}{\ensuremath{\mathtt{qcase}}}
\newcommand{\textmatch}{\ensuremath{\mathtt{match}}}
\newcommand{\textletrec}{\ensuremath{\mathtt{letrec}}}
\newcommand{\textshape}{\ensuremath{\mathtt{shape}}}
\newcommand{\zket}{\ket 0}
\newcommand{\oket}{\ket 1}
\newcommand{\qcasebase}[2]{
  \textqcase\,#1\left\{\! \begin{aligned} #2 \end{aligned}\right\}
}
\newcommand{\qcase}[3]{
  \qcasebase {#1} {\zket \to #2\,, \oket \to #3}
}
\newcommand{\qcasesplit}[3]{
  \qcasebase {#1} { 
    \begin{aligned} 
      \zket &\to #2 \\ \oket &\to #3 
    \end{aligned}
  }
}
\newcommand{\qcasedefault}[1]{\ensuremath{\qcase{#1}{t_0}{t_1}}}
\newcommand{\match}[2]{
  \textmatch\,#1\left\{\! \begin{aligned} #2 \end{aligned}\right\}
}
\newcommand{\matchdefault}[1]{\textmatch\,#1 \left\{
    c_1(\overrightarrow{x_1}) \to t_1\,, \dots,
    c_n(\overrightarrow{x_n}) \to t_n
\right\}}
\newcommand{\matchdefaultcollapsedany}[2]{
  \textmatch_{1 \leq i \leq n}\,#1\,\{
    c_i(\overrightarrow{x_i}) \to #2_i
  \}
}
\newcommand{\matchdefaultcollapsed}[1]{\matchdefaultcollapsedany{#1}{t}}
\newcommand{\lbd}[2]{\lambda #1 .#2}
\newcommand{\letrec}[3]{\textletrec\,#1\,#2=#3}
\newcommand{\shape}[1]{\mathtt{shape}(#1)}

%Terms sets and properties
\newcommand{\can}{\ensuremath{\mathtt{CAN}}}
\newcommand{\equivcontext}{\ensuremath{C_\equiv}}
\newcommand{\fv}[1]{\mathrm{FV}(#1)}
\newcommand{\varset}{\mathtt{Var}}
\newcommand{\valueset}{\mathrm{Value}}
\newcommand{\subterm}{\vartriangleleft}
\newcommand{\subtermeq}{\trianglelefteq}

%Typing
\newcommand{\vartyping}[2]{#1 : #2}
\newcommand{\typing}[4]{
  \ifx&#1&\varnothing \else #1\fi; 
  \ifx&#2&\varnothing \else #2\fi 
  \vdash {#3} : #4
}
\newcommand{\closedtyping}[2]{\typing{}{}{#1}{#2}}
\newcommand{\shapemarker}[1]{[#1]}
\newcommand{\typetermset}{\termset_{\vdash}}
\newcommand{\typedom}[1]{\mathrm{dom}(#1)}
%Orthogonality
\newcommand{\pitwo}{\Pi_2^0}
\newcommand{\portho}{\mathtt{ORTHO}}
\newcommand{\uhalt}{\mathtt{UHalt}}
\newcommand\kron[2]{\delta_{#1, #2}}

%Unitarity
\newcommand\innerprod[2]{\langle #1, #2 \rangle}
\newcommand\innerprodblk{\innerprod \cdot \cdot \qvar}

%Substitutions
\newcommand{\subarrow}{\to}
\newcommand{\subsupp}[1]{\mathrm{Supp}(#1)}

%Reduction
\newcommand{\reduces}{\rightsquigarrow}
\newcommand{\sreduces}{\reduces^*}
\newcommand{\evalcontext}{\ensuremath{E}}
\newcommand{\zoreduces}{\ensuremath{\reduces^?}}%Zero or one reduction
\newcommand{\mostreduces}[1]{\ensuremath{\reduces^{\leq #1}}}

%Circuit compilation
\newcommand{\fbqpscript}{\mathtt{QList}}
\newcommand{\circuittype}{\ensuremath{\mathtt{CircuitTerms}}}
\newcommand{\hyrqlbound}[1]{\ensuremath{\mathtt{Hyrql}({#1})}}
\newcommand{\hyrqlpoly}{\hyrqlbound{\mathrm{Poly}}_\fbqpscript}
\newcommand{\binfunc}{\set{0,1}^* \to \set{0,1}^*}
\newcommand{\boundedrec}{\ensuremath{\mathtt{Faithful}}}
\newcommand{\shapeset}[1]{\mathtt{S}(#1)}
\newcommand{\famcirc}{\mathtt{C}(t)}
\newcommand{\bqwalk}{\mathtt{bqwalk}}
\newcommand{\timeset}[1]{\mathtt{Time}(#1)}

%Misc
\newcommand{\ie}{\text{i.e., }}
\newcommand{\eg}{\text{e.g., }}
\newcommand{\gmid}{\,\mid\,}
\newcommand{\size}[1]{|#1|}
\newcommand{\QS}{\ensuremath{\mathtt{QS}}}
\newcommand{\hyrql}{{{\tt{Hyrql}}}}
\newcommand{\spm}{{{\tt{Spm}}}}
\newcommand{\mcal}{\mathcal}
\newcommand{\seq}[1]{\overrightarrow{#1}}
\newcommand{\fbqp}{\ensuremath{\mathtt{FBQP}}}
\newcommand{\bigo}[1]{\mathcal O({#1})}
\newcommand{\mbN}{\mathbb N}
\newcommand{\mtf}{\mathtt f}
\newcommand{\mcC}{\mathcal C}
\newcommand{\computes}[1]{\llbracket #1 \rrbracket}
%Proofs
\newcommand{\qt}[2]{\theta_{#1}(#2)}
\newcommand{\kroneq}[2]{\nu_{#1, #2}}

% Delimiters
\DeclarePairedDelimiter\bbracket{\llbracket}{\rrbracket}
\DeclarePairedDelimiter\floor{\lfloor}{\rfloor}
\renewcommand{\subsectionautorefname}{Subsection}

\newcommand{\Emmanuel}[1]{\todo{#1 (Emmanuel)}}
\newcommand{\Thomas}[1]{\todo{(Thomas) #1}}
\newcommand{\Romain}[1]{\todo{#1 (Romain)}}


\newtheorem{theorem}{Theorem}
\newtheorem{lemma}[theorem]{Lemma}
\newtheorem{proposition}[theorem]{Proposition}
\newtheorem{corollary}[theorem]{Corollary}
\newtheorem{definition}[theorem]{Definition}
\newtheorem{example}[theorem]{Example}
\newtheorem{remark}[theorem]{Remark}

